% !TEX root = ../../../main.tex
% 来源：中学数学实验教材·第三册（高中数学）
% 章节：任意角的三角函数 / 任意角的三角函数
% 原始文件：6.tex，行 948-982
% 类型：tikz
% ---
\begin{tikzpicture}[>=latex]
\begin{scope}
\draw[->] (-1.5,0)--(1.5,0)node[right]{$x$};
\draw[->] (0,-1.5)--(0,1.5)node[right]{$y$};
\node at (-.25,-.25){$O$};
\draw (0,0) circle(1);
\foreach \x/\xcorr in {+/{.5,.5}, -/{-.5,.5}, -/{-.5,-.5}, +/{.5,-.5}}
{
    \node at (\xcorr) {$\x$};
}
\node at (0,-2){$\cos\alpha$和$\sec\alpha$};
\end{scope}
\begin{scope}[xshift=4cm]
    \draw[->] (-1.5,0)--(1.5,0)node[right]{$x$};
    \draw[->] (0,-1.5)--(0,1.5)node[right]{$y$};
    \node at (-.25,-.25){$O$};
    \draw (0,0) circle(1);
    \foreach \x/\xcorr in {+/{.5,.5}, +/{-.5,.5}, -/{-.5,-.5}, -/{.5,-.5}}
    {
        \node at (\xcorr) {$\x$};
    }
    \node at (0,-2){$\sin\alpha$和$\csc\alpha$};
    \end{scope}
    \begin{scope}[xshift=8cm]
\draw[->] (-1.5,0)--(1.5,0)node[right]{$x$};
\draw[->] (0,-1.5)--(0,1.5)node[right]{$y$};
\node at (-.25,-.25){$O$};
\draw (0,0) circle(1);
\foreach \x/\xcorr in {+/{.5,.5}, -/{-.5,.5}, +/{-.5,-.5}, -/{.5,-.5}}
{
    \node at (\xcorr) {$\x$};
}
\node at (0,-2){$\tan\alpha$和$\cot\alpha$};
\end{scope}
\end{tikzpicture}
